% 第6章 Work 2：LLM KV缓存

\FloatBarrier

\chapter{大语言模型KV缓存的有机光电存内计算映射}
\label{ch:work2_scope}

\FloatBarrier

\section{研究动机与物理匹配}
\label{sec:work2_motivation}
\label{sec:work2_physics_match}

本章的出发点很简单：KV 缓存不是模型权重。它随上下文长度线性增长，在解码阶段被反复读取，并在会话结束后释放，因此主要瓶颈是存储容量、数据搬运和会话级保持，而不是长期非易失保存。其容量可写为
\begin{equation}
M_{\text{KV}} = 2 \times L \times H \times D_h \times S \times P,
\label{eq:kv_memory}
\end{equation}
其中 $L$ 为 Transformer 层数，$H$ 为注意力头数，$D_h$ 为每头维度，$S$ 为序列长度，$P$ 为单元素字节数。以 LLaMA-2-7B 为例，$L=32$、$H=32$、$D_h=128$，在 $S=4096$ 且 FP16 存储时，单批次 KV 缓存约为 2~GiB；当上下文扩展至 32k 或批量解码时，缓存容量可超过模型权重本身。

现有数字路线分别压缩、驱逐或搬运 KV 缓存，但没有消除这一存储--寿命错配。KIVI~\cite{liu2024kivi} 等低比特量化需要 scale/zero-point 与反量化逻辑；StreamingLLM~\cite{xiao2024streamingllm}、H$_2$O~\cite{zhang2023h2o} 依赖可预测的注意力稀疏性；InfiniGen~\cite{lee2024infinigen}、vLLM~\cite{kwon2023efficient} 和 FlashAttention~\cite{dao2022flashattention} 优化搬运路径，却仍需管理持续增长的缓存状态。由此，本章不主张用有机存储替代通用 DRAM，而是考察它能否作为选择性、会话级 KV 存内计算层。

解码阶段的注意力计算为
\begin{equation}
s_{ti} = \frac{q_t \cdot k_i^T}{\sqrt{D_h}}, \quad \alpha_{ti} = \frac{\exp(s_{ti})}{\sum_{j=1}^{S} \exp(s_{tj})},
\label{eq:attention_scores}
\end{equation}
\begin{equation}
o_t = \sum_{i=1}^{S} \alpha_{ti} v_i .
\label{eq:attention_output_detail}
\end{equation}
$q_tK^T$ 和 $\alpha_tV$ 本质上都是矩阵--向量运算，天然适合交叉阵列的欧姆定律--基尔霍夫定律计算。KV 缓存的访问模式也与权重推理不同：K/V 条目在 prefill 或新 token 生成时写入，随后在会话内被反复读取，形成低频写入、高频读取、会话级释放的 WORM（write-once-read-many）负载。写延迟在这种负载下被摊薄，而读能耗和数据搬运才是主导项，这使得存内计算从通用加速器问题转化为一个更具体的存储--负载匹配问题。

不同存储技术在保持时间、耐久性和刷新机制上与 KV 缓存需求的匹配程度不同。令
\begin{equation}
\eta_{\tau} = \frac{\tau_{\text{device}}}{\tau_{\text{session}}}, \quad \eta_{E} = \frac{E_{\text{max}}}{N_{\text{write}}},
\label{eq:match_metrics}
\end{equation}
分别表示保持时间裕度和耐久性裕度。理想器件并不是在两个维度上无限大，而是在满足会话需求的同时避免为多年保持或百万级写入循环支付不必要代价。表~\ref{tab:kv_tech_match} 给出了典型技术的定性比较。

\begin{table}[!htbp]
\centering
\caption{不同存储技术与 KV 缓存会话级 WORM 负载的匹配度。KV 缓存需要低频写入、高频读取和秒至分钟级保持；有机 OEC-RAM 的有限保持时间与低耐久性在该负载下从缺陷转化为可利用的物理匹配。}
\label{tab:kv_tech_match}
\begin{tabular}{@{}lcccc@{}}
\toprule
技术 & 保持时间 & 耐久性 & 面积效率 & 光学刷新 \\
\midrule
SRAM & ms 级（易失） & $>10^{15}$ & 低（6T/8T 单元） & 不支持 \\
氧化物 RRAM（HfO$_2$） & 年（$>10^8$ s） & $>10^{6}$ & 中 & 不支持 \\
PCM（Ge$_2$Sb$_2$Te$_5$） & 年（$>10^8$ s） & $>10^{8}$ & 中 & 不支持 \\
FeFET & 月--年 & $>10^{5}$ & 高 & 不支持 \\
有机 OEC-RAM & $\sim 10^3$ s & $10^3$--$10^4$ & 高 & 支持 \\
\bottomrule
\end{tabular}
\end{table}

SRAM 速度快但面积和静态功耗随容量线性增长；GB 级 KV 缓存若全部放在 SRAM 中，端侧面积和漏电成本难以接受。氧化物 RRAM、PCM 和 FeFET 具备非易失性，但其年级保持时间和高耐久性对于秒至分钟级会话是过度规格，通常伴随更高编程电压、更复杂 forming 或更严格写入校验。相比之下，有机 OEC-RAM 的 $\sim 10^3$~s 保持时间与会话长度处于同一数量级，$10^3$--$10^4$ 次循环虽低于无机非易失存储，却足以覆盖 KV 缓存的低频写入需求。一次会话内的 K/V 元素数量很大，但阵列编程以向量和块为单位并行完成，器件层面的重复编程次数远低于元素计数。

有机光电器件的光敏性在视觉推理中可能是噪声源，在 KV 缓存中则提供了低耐久成本的数据维持方式。光学刷新可用光脉冲把电导态推回目标区间，不需要像电学重写那样消耗完整编程循环；通过波导或局部照明，还可优先刷新近期 token 或高注意力 token 所在子阵列。本文因此把 Work~2 的中心问题定义为：有机光电 CIM 是否能在不声称替代通用 DRAM 的前提下，作为会话级 KV 缓存的局部存储--计算层，在选择性层范围内获得可接受困惑度和更好的物理代价。

\FloatBarrier

\section{问题定义与评估框架}
\label{sec:work2_problem}
\label{sec:work2_models}

设边缘 LLM 推理由预填充阶段 $\mathcal{P}$ 和解码阶段 $\mathcal{D}$ 组成。预填充阶段逐层产生
\begin{equation}
\mathbf{K}_l = \mathbf{W}_l^K \mathbf{x}, \quad \mathbf{V}_l = \mathbf{W}_l^V \mathbf{x}, \quad l = 1, \dots, L,
\label{eq:kv_compute}
\end{equation}
解码阶段第 $t$ 步输出
\begin{equation}
\mathbf{o}_t = \sum_{i=1}^{S+t-1} \alpha_{ti} \mathbf{v}_i, \quad \alpha_{ti} = \frac{\exp(q_t k_i^T / \sqrt{D_h})}{\sum_j \exp(q_t k_j^T / \sqrt{D_h})}.
\label{eq:attention_detail}
\end{equation}
当 K/V 条目被映射为交叉阵列中的电导 $G_{ij}$ 后，读出电流为
\begin{equation}
I_j = \sum_i W_{ij} V_i, \quad W_{ij} \propto G_{ij}.
\label{eq:crossbar_mvm}
\end{equation}
ADC 将电流转换回数字域，softmax 与后续控制逻辑仍保留在数字端。该混合划分与前文视觉模型一致：密集线性运算交给模拟阵列，动态归一化和控制流留在数字域。

本文围绕三项假设组织 Work~2。保持时间匹配假设认为，有机 OEC-RAM 的 $\tau\sim10^3$~s 与边缘 LLM 会话长度 $\tau_{\text{session}}\in[10^1,10^3]$~s 同阶，因此无需为长期非易失性支付额外物理代价，也不会在典型会话内过早丢失缓存内容。
\begin{hypothesis}[保持时间匹配假设]
若 $\tau_{\text{device}}/\tau_{\text{session}}\in[1,100]$，则有机 KV 缓存可覆盖会话周期内的自然衰减，并避免多年保持带来的写能耗过度规格。
\label{hyp:retention}
\end{hypothesis}
耐久性适配假设认为，KV 缓存的 WORM 模式把器件重复写入次数限制在 $E_{\text{max}}$ 以下，光学刷新进一步降低了电学重写频率。
\begin{hypothesis}[耐久性适配假设]
若 $\eta_E=E_{\text{max}}/N_{\text{write}}\ge10$，则系统可在完整会话内维持电导精度而不耗尽有机器件的耐久预算。
\label{hyp:endurance}
\end{hypothesis}
模拟精度充分假设则不能只靠器件参数推出，必须通过行为级实验检验。
\begin{hypothesis}[模拟精度充分假设]
存在有限模拟态数 $N\ge N_{\text{min}}$ 和噪声上限 $\sigma\le\sigma_{\text{max}}$，使有机 CIM KV 缓存相对 FP16 基线的困惑度和下游任务退化低于应用阈值 $\delta$。
\label{hyp:precision}
\end{hypothesis}
三项假设对应器件物理、系统访问模式和算法精度三个层次。前两项由有机 OEC-RAM 的保持与耐久数据提供先验支持，第三项需要通过噪声、量化和层选择实验确定操作窗口。

主模型选择 TinyLlama-1.1B~\cite{zhang2024tinyllama}。该模型包含 22 层 Transformer、隐藏维度 2048、32 个注意力头、每头维度 64；在 4k 上下文、FP16 KV 缓存下，缓存约为 704~MB，适合作为 1--3B 边缘模型的代表。LLaMA-2-7B~\cite{touvron2023llama2} 仅作为扩展验证，用于观察 GQA/MQA 结构和更大隐藏维度对阵列尺寸、ADC 和缓存容量的压力。当前范围不纳入 $>$7B 模型或 $>$32k 上下文，因为这些情形需要多片划分、3D 堆叠或片外卸载，已经超出器件物理匹配本身。

评估以 WikiText 困惑度为主，并在 LongBench~\cite{bai2023longbench}、GSM8K~\cite{cobbe2021training} 和 MT-Bench~\cite{zheng2023judging} 子集上检查长上下文问答、数学推理和多轮对话的下游影响。系统侧报告每 token 解码能耗、KV 缓存面积和有效保持时间。所有有机 CIM 结果应使用固定随机种子、归档配置文件、原始日志和可复现实验命令；缺少 git commit、脚本、数据集、评估协议或 checkpoint hash 的阶段性结果只能用于路线筛选，不升级为论文锁定结论。

\begin{table}[!htbp]
\centering
\caption{对比基线及公平比较原则。}
\label{tab:baselines}
\small
\begin{tabularx}{\textwidth}{@{}l>{\raggedright\arraybackslash}X>{\raggedright\arraybackslash}X@{}}
\toprule
基线 & 技术描述 & 公平比较原则 \\
\midrule
FP16 & 标准数字 DRAM/SRAM 存储，半精度浮点 & 精度上界；本文只讨论以更低物理代价接近 FP16 \\
KIVI-2bit & 2-bit 非对称 KV 量化~\cite{liu2024kivi} & 在等 footprint 下比较困惑度与反量化成本 \\
KIVI-4bit & 4-bit 非对称 KV 量化 & 作为精度--压缩中间点 \\
SRAM-CIM & 7~nm SRAM 交叉阵列~\cite{hastily2023} & 计入面积和待机功耗，不只比较读延迟 \\
氧化物 RRAM-CIM & HfO$_2$ RRAM 交叉阵列~\cite{huang2024retransformer} & 比较写能耗、保持时间、耐久性和模拟态有效数 \\
\bottomrule
\end{tabularx}
\end{table}

\FloatBarrier

\section{实验设计}
\label{sec:work2_experiments}

Work~2 的实验按从精度到系统、从静态到时序的顺序展开。第一组实验测量注意力对模拟 KV 存储噪声的容忍度；第二组把该容忍度转化为阵列尺寸、ADC 位宽和能耗面积约束；第三组在等 footprint 下比较有机模拟存储与数字 KV 量化；第四组加入保持漂移和刷新策略。四组实验共享 Work~1 的器件参数文件和噪声接口，以避免在不同章节之间改变物理假设。

\noindent\textbf{实验一：注意力噪声鲁棒性。}\label{sec:exp1_noise}
在 TinyLlama-1.1B 的 HuggingFace 实现中，噪声钩子插入各层自注意力模块的 $k_{\text{proj}}$ 和 $v_{\text{proj}}$ 输出之后。浮点 K/V 张量先映射为有限电导码本，再叠加固定实例级 D2D 偏置和每次前向传播独立采样的 C2C 扰动：
\begin{equation}
G_{\text{noisy}} = G + \mathcal{N}(0, \sigma_{\text{c2c}}^2) + \mathcal{N}(0, \sigma_{\text{d2d}}^2).
\label{eq:noise_model}
\end{equation}
扫描空间见表~\ref{tab:exp1_grid}。除困惑度外，实验计算注意力 top-$k$ 保持率
\begin{equation}
\rho_{\text{rank}} = \frac{|\text{top-}k(A_{\text{FP16}}) \cap \text{top-}k(A_{\text{organic}})|}{k},
\label{eq:rank_preservation}
\end{equation}
用于区分数值误差和注意力排序结构改变。若 $\sigma_{\text{c2c}}$ 或 $\sigma_{\text{d2d}}$ 稍增即导致困惑度陡升，则本文将把声称限定为有界噪声下的 graceful degradation，而不使用无损注意力表述。

\begin{table}[!htbp]
\centering
\caption{实验一参数扫描空间。}
\label{tab:exp1_grid}
\begin{tabular}{@{}llc@{}}
\toprule
维度 & 取值 & 说明 \\
\midrule
模拟态数 $N$ & 2, 3, 4, 5, 6 & 每器件电导离散级数 \\
C2C 噪声 $\sigma_{\text{c2c}}$ & 0.00, 0.01, 0.02, 0.05, 0.10 & 相对标准差 \\
D2D 噪声 $\sigma_{\text{d2d}}$ & 0.00, 0.05 & 固定实例偏置 \\
注入目标 & K-only, V-only, K+V & 分离 Key 与 Value 敏感性 \\
\bottomrule
\end{tabular}
\end{table}

\noindent\textbf{实验二：架构探索与帕累托分析。}\label{sec:exp2_arch}
第二组实验扩展行为级模拟器，建立 KV 缓存交叉阵列模块。每个配置执行单 token、单层、单头的注意力计算，并输出延迟、能耗、面积和 SNR。模拟 MAC 能耗按 Work~1 的一阶模型继承，ADC 采用 Walden 优值随位宽指数缩放，写入和光学刷新分别按 pJ/单元估算。表~\ref{tab:exp2_grid} 的 $108$ 组配置用于确定精度约束帕累托前沿：只有满足实验一给出的 $\sigma_{\text{max}}$ 或 SNR 下限的配置才进入候选集合。

\begin{table}[!htbp]
\centering
\caption{实验二参数扫描空间。}
\label{tab:exp2_grid}
\begin{tabular}{@{}ll@{}}
\toprule
维度 & 取值 \\
\midrule
阵列尺寸 & $64\times64$, $128\times128$, $256\times256$, $512\times512$ \\
ADC 位数 & 4, 6, 8 \\
编程精度 & 1\%, 5\%, 10\%（$\sigma_{\text{prog}}/G_{\text{range}}$） \\
模拟态数 & 4, 8, 16 \\
\bottomrule
\end{tabular}
\end{table}

\noindent\textbf{实验三：压缩与模拟存储协同设计。}\label{sec:exp3_compression}
第三组实验在等效存储密度 $\rho$ 下比较数字量化和有机模拟存储。数字侧采用 KIVI 式非对称量化；有机侧比较均匀电报码本和基于校准集 K-means 的学习码本。对于数字方案，$\rho$ 等于量化位数；对于有机方案，$\rho=\log_2N$。比较指标包括注意力输出余弦相似度、困惑度退化、任务指标和每 token 写能耗。表~\ref{tab:exp3_grid} 的扫描用于判断非均匀码本是否能弥补模拟噪声 floor，并定位 2--4 bit 等效区间中的优势和劣势。

\begin{table}[!htbp]
\centering
\caption{实验三参数扫描空间。}
\label{tab:exp3_grid}
\begin{tabular}{@{}ll@{}}
\toprule
维度 & 取值 \\
\midrule
数字量化位数 & 2, 3, 4 bit \\
有机模拟态数 & 4, 8, 16 \\
码本类型 & 均匀, 学习（K-means） \\
器件噪声 & $\sigma_{\text{c2c}}\in\{0.00,0.01,0.02\}$, $\sigma_{\text{d2d}}\in\{0.00,0.02,0.04\}$ \\
\bottomrule
\end{tabular}
\end{table}

\noindent\textbf{实验四：时间动态与刷新策略。}\label{sec:exp4_temporal}
第四组实验把保持漂移写为
\begin{equation}
G(t) = G_{\min} + (G_{\text{programmed}} - G_{\min}) \cdot \exp(-t / \tau),
\label{eq:retention_decay}
\end{equation}
并在多轮对话中比较无刷新、周期刷新和驱逐策略。评估分三层进行：单头级别测量注意力分数相对误差，模型级别测量困惑度累积，任务级别测量多轮对话胜率和总能耗。若 $\tau\ge100$~s 且刷新间隔 $\le10$~s 时仍能维持应用精度，则假设~\ref{hyp:retention} 获得时间动态层面的支持；若频繁刷新导致能耗或耐久预算失控，则有机 KV 缓存只能作为短会话候选。

\begin{table}[!htbp]
\centering
\caption{实验四参数扫描空间。}
\label{tab:exp4_grid}
\begin{tabular}{@{}ll@{}}
\toprule
维度 & 取值 \\
\midrule
保持时间常数 $\tau$ & 1 s, 10 s, 100 s, 1000 s \\
刷新间隔 & 无, 1 s, 10 s, 100 s \\
驱逐策略 & FIFO, LRU, 注意力分数阈值 \\
存储层级 & 单层有机 CIM, 有机 CIM L1 $\rightarrow$ SRAM L2 \\
\bottomrule
\end{tabular}
\end{table}

\FloatBarrier

\section{Remote 107阶段性结果}
\label{sec:work2_results}

Remote~107 分支在 2026 年 4--5 月完成了选择性 KV 缓存验证。需要先划清证据边界：旧 TinyLlama/候选 lane 只解释路线筛选过程，不能与当前锁定数据混表；当前可作为论文依据的是 selective-KV claim-lock packet，其协议为 Pythia-410M、WikiText-2 raw-v1 test、context length 512、non-overlapping stride 512、batch size 1，数字参考 PPL=23.30，并已归档 41/41 个逐行 JSON artifact 及对应 manifest 元数据。该锁定包回答的核心问题是层范围选择：终端层 KV 可以继续推进，全层模拟 KV 应作为负面对照。

历史全层 8-bit 零噪声候选结果曾显示相对旧 15.68 基线退化 11.5\%，见表~\ref{tab:work2_baseline}。这一结果不再承担当前数值主结论，但保留为路线筛选证据：即使不加入器件噪声，全层模拟化也不是稳健主线。

\begin{table}[!htbp]
\centering
\caption{Remote~107 旧候选 lane 给出的全层淘汰信号（seed=42）。该表使用旧 15.68 数字基线，只用于说明全层模拟 KV 在早期筛选中已不稳健；最终数值结论以后续 claim-lock packet 为准。}
\label{tab:work2_baseline}
\begin{tabular}{@{}lccc@{}}
\toprule
配置 & PPL & 相对基线 & 判定 \\
\midrule
数字基线（FP16） & 15.68 & 1.000x & 参考 \\
8-bit 全层零噪声 & 17.48 & 1.115x & 淘汰 \\
\bottomrule
\end{tabular}
\end{table}

选择性终端层显著缓和了这一问题。last1（layer~23）在零噪声下 PPL=15.82，相对基线仅退化 0.9\%；真实噪声下 PPL=16.72，退化 6.6\%，仍处于 10\% 门限内，见表~\ref{tab:work2_selective}。last2 也保持可观察价值，但一致性略弱于 last1。

\begin{table}[!htbp]
\centering
\caption{选择性终端层模拟 KV 缓存探测（8-bit, seed=42）。}
\label{tab:work2_selective}
\begin{tabular}{@{}lccc@{}}
\toprule
层范围 & 噪声条件 & PPL & 相对基线 \\
\midrule
last1（layer~23） & 零噪声 & 15.82 & 1.009x \\
last1（layer~23） & 真实噪声 & 16.72 & 1.066x \\
\bottomrule
\end{tabular}
\end{table}

\begin{figure}[!htbp]
\centering
\includegraphics[width=0.90\textwidth]{paper2_fig_107_kv_candidate_summary}
\caption{Remote~107 选择性 KV 缓存候选路线的阶段性两联图。A 比较 last1/last2/last4/all24 在不同 D2D 候选条件下的困惑度，B 显示全层配置对 D2D 噪声的敏感性；数字基线 PPL=15.68 仅作为参考。该图为审计草稿，尚未满足完整元数据锁定要求。}
\label{fig:paper2-107-kv-candidate-summary}
\end{figure}

图~\ref{fig:paper2-107-kv-candidate-summary} 保留早期候选 lane 的审计草稿，用于说明为何全层模拟化被降级为负面对照。当前可作为论文锁定引用的是随后完成的 selective-KV claim-lock packet，其协议为 Pythia-410M、WikiText-2 raw-v1 test、context length 512、non-overlapping stride 512、batch size 1，数字参考为 PPL=23.30，并且逐行命令、checkpoint hash、seed 语义、corrected-noise 代码路径与 41 个原始 JSON artifact 已完成归档。表~\ref{tab:work2_107_clean_scope} 和图~\ref{fig:paper2-107-clean-protocol} 单独报告该锁定包，不与旧 15.68 基线或本地 cache-path probe 混合。该协议下排序清晰：$\sigma_{\text{d2d}}=0.02$ 时 last1 为 $18.781\pm0.023$，last2 为 $18.594\pm0.014$，last4 为 $20.402\pm0.069$，全层 all24 为 $25.985\pm0.148$；$\sigma_{\text{d2d}}=0.05$ 时 last1 与 last2 仍分别为 $18.963\pm0.035$ 和 $19.018\pm0.025$，last4 升至 $21.535\pm0.051$，全层则恶化到 $59.055\pm1.610$。这组锁定数据支持“选择性终端层可继续推进、全层模拟化只保留为负面对照”的结论。

\begin{table}[!htbp]
\centering
\caption{Remote~107 selective-KV claim-lock packet 下的层范围排序。Pythia-410M 在 WikiText-2 raw-v1 test、context length 512、non-overlapping stride 512、batch size 1 下的数字参考为 PPL=23.30；last1/last2 在两档 D2D 下保持稳定，而 all24 在 D2D=0.05 时失控，因此本章把终端层选择性缓存作为主路线、全层模拟缓存作为负面对照。}
\label{tab:work2_107_clean_scope}
\begin{tabular}{@{}lcccc@{}}
\toprule
层范围 & 评估 D2D & $n$ & PPL & 相对 23.30 基线 \\
\midrule
last1 & 0.02 & 5 & $18.781\pm0.023$ & 0.806x \\
last1 & 0.05 & 5 & $18.963\pm0.035$ & 0.814x \\
last2 & 0.02 & 5 & $18.594\pm0.014$ & 0.798x \\
last2 & 0.05 & 5 & $19.018\pm0.025$ & 0.816x \\
last4 & 0.02 & 5 & $20.402\pm0.069$ & 0.876x \\
last4 & 0.05 & 5 & $21.535\pm0.051$ & 0.924x \\
all24 & 0.02 & 5 & $25.985\pm0.148$ & 1.115x \\
all24 & 0.05 & 5 & $59.055\pm1.610$ & 2.535x \\
\bottomrule
\end{tabular}
\end{table}

\begin{figure}[!htbp]
\centering
\includegraphics[width=0.90\textwidth]{paper2_fig_107_clean_protocol_summary_20260511}
\caption{Remote~107 selective-KV claim-lock packet 的选择性 KV 缓存结果。A 显示 last1/last2/last4/all24 随 D2D 失配增加的困惑度变化，B 直接比较选择性终端层与全层模拟 KV；协议为 Pythia-410M、WikiText-2 raw-v1 test、context length 512、non-overlapping stride 512、batch size 1，数字参考 PPL=23.30，并已归档 41/41 个逐行 JSON artifact。图中结论只锁定层范围排序和全层负面对照，不声称能耗、耐久或长上下文部署已经闭环。}
\label{fig:paper2-107-clean-protocol}
\end{figure}

Remote~107 随后的 R3 自适应调度 lane 提供了更细的工程线索，但当前仍应保持为验证性证据，而不是第二张主 claim 表。对冻结非目标参数的 410M 主线而言，fixed、cosine 和 layer-wise 三种调度的训练后 PPL 分别为 18.75、18.31 和 18.36，而 reverse layer-wise 恶化到 21.30，说明把更强噪声放在浅层并不能形成有效正则。对更大模型而言，这一差距明显缩小：Pythia-2.8B 为 12.68/12.68/12.70，Pythia-6.9B 为 11.40/11.38/11.43，提示主路线在规模增大后对调度形状不再高度敏感。

\begin{figure}[!htbp]
\centering
\includegraphics[width=0.90\textwidth]{paper2_fig_remote107_r3_adaptive_ppl_20260513}
\caption{Remote~107 R3 自适应噪声调度比较。对于冻结非目标参数的 410M 主线，cosine decay 达到最低 PPL=18.31，略优于 layer-wise 的 18.36 和 fixed 的 18.75；reverse layer-wise 退化到 21.30，作为负面对照。对 2.8B 和 6.9B，fixed、cosine 和 layer-wise 之间只保留小幅差异。}
\label{fig:paper2-107-r3-adaptive}
\end{figure}

真正出现结构性收益的是扩大模拟覆盖范围时的调度选择。Pythia-2.8B 的联合消融表明，last1 fixed 为 12.68，last1 cosine 同样为 12.68，last2 fixed 升至 13.78，而 last2+cosine 进一步降到 12.56，成为当前 p28b 路线中的最佳结果；last4 fixed 则回升到 14.12。也就是说，R3 没有推翻“last1 是默认主路线”的结论，而是补充了一条更细的设计规则：若希望把模拟 KV 的覆盖范围从单一终端层扩展到 last2，以换取更高的面积或能效收益，应优先采用自适应 schedule，而不是简单沿用 fixed 训练。

\begin{figure}[!htbp]
\centering
\includegraphics[width=0.82\textwidth]{paper2_fig_remote107_r3_p28b_ablation_20260513}
\caption{Pythia-2.8B 层范围与调度联合消融。last2+cosine 达到最佳 PPL=12.56，优于 last2+fixed 的 13.78，也略优于 last1+fixed 的 12.68。该结果支持把自适应调度作为扩大模拟覆盖范围时的工程补偿手段。}
\label{fig:paper2-107-r3-p28b}
\end{figure}

作为与困惑度不同的离线诊断，图~\ref{fig:paper2-kv-offline-bitwidth} 直接读取 Pythia-410M 的真实 KV 张量，经过量化、固定 D2D 和重复 C2C 读出后计算相对重构均方误差。该结果显示 4-bit 到 6/8-bit 存在明显误差下降，16-bit 接近当前噪声底；last2 与 last4 的 8-bit 相对重构误差约为 $1.2\times10^{-4}$，与 last1 处于同一数量级。这一图只说明张量存储--读出链路的数值误差结构，不能替代端到端 PPL 评估。

\begin{figure}[!htbp]
\centering
\includegraphics[width=0.90\textwidth]{paper2_fig_w2_kv_offline_bitwidth_scan_20260511}
\caption{真实 KV 张量离线重构的位宽诊断。A 显示 last1/last2/last4/all24 在 4、6、8、16-bit 下的相对 KV 重构均方误差，B 比较 8-bit 与 16-bit 的层范围差异。该实验使用真实 Pythia-410M KV 张量和模拟读出噪声，但没有接入注意力前向传播，因此只作为数值诊断而非困惑度证据。}
\label{fig:paper2-kv-offline-bitwidth}
\end{figure}

图~\ref{fig:paper2-kv-retention-refresh} 进一步把离线诊断扩展到时间保持与刷新。该实验使用同一批真实 Pythia-410M KV 张量，在 8-bit、$\sigma_{\mathrm{D2D}}=0.005$、$\sigma_{\mathrm{C2C}}=0.002$ 下扫描会话跨度、保持时间常数和刷新间隔。无刷新且会话跨度远大于 $\tau$ 时，相对重构误差可接近 1，说明缓存态已经显著衰减；当 $\tau=1000$~s 且每 10~s 刷新时，last1/last2/last4/all24 在 1000~s 跨度下的误差回落到约 $10^{-4}$--$10^{-3}$。该图为假设~\ref{hyp:retention} 的数值压力测试提供了方向，但仍没有接入端到端注意力前向传播，也没有核算刷新能耗和耐久预算。

\begin{figure}[!htbp]
\centering
\includegraphics[width=0.90\textwidth]{paper2_fig_w2_kv_retention_refresh_scan_20260511}
\caption{真实 KV 张量离线重构的保持--刷新诊断。A 显示 1000~s 会话跨度且无刷新时，相对 KV 重构误差随保持时间常数 $\tau$ 改善；B 显示在 $\tau=1000$~s 下，不同刷新间隔对 1000~s 会话跨度的误差恢复作用。该实验只评估存储--读出链路的时间动态，不构成端到端困惑度或能耗结论。}
\label{fig:paper2-kv-retention-refresh}
\end{figure}

表~\ref{tab:work2_alllayer_sweep} 的 6-bit 全层扫描给出了更强的压力证据：C2C 和 D2D 任一轴增大都会快速推高困惑度，尤其 D2D=0.10 时已完全失去实用意义。500 步 HAT 预热可把全层 PPL 从 579.52 降至 142.27，但仍远未通过门限，因此 HAT 资源应集中在选择性终端层，而非全层救援。

\begin{table}[!htbp]
\centering
\caption{6-bit 全层噪声敏感性扫描（seed=42）。}
\label{tab:work2_alllayer_sweep}
\begin{tabular}{@{}lcccccc@{}}
\toprule
$\sigma_{\text{c2c}}$ & 0.0 & 0.005 & 0.01 & 0.015 & 0.02 & 0.03 \\
PPL & 32.41 & 39.86 & 56.73 & 77.36 & 100.64 & 160.34 \\
\midrule
$\sigma_{\text{d2d}}$ & 0.0 & 0.01 & 0.02 & 0.05 & 0.10 & \\
PPL & 32.41 & 60.30 & 106.48 & 470.92 & 7635.49 & \\
\bottomrule
\end{tabular}
\end{table}

Fresh-D2D 交叉实例验证进一步支持选择性路线。训练时固定 D2D=0.02，评估时在 5 个独立 D2D 种子上重复。表~\ref{tab:work2_fresh_d2d} 显示，last1 在 eval D2D=0.02--0.05 范围内仅从 18.42 增至 18.60，标准差始终低于 0.03；last2 从 18.71 增至 19.21，仍表现出低方差但均值略高；全层则从 26.05 增至 66.84，且跨实例方差快速扩大。当前路线排序因此是 last1 为主、last2 为备选、全层仅作压力测试。高 D2D 训练对全层有一定收益，例如 D2D=0.04 训练在 eval 0.05 时优于 D2D=0.02 训练，但对 last1 没有明显必要。

\begin{table}[!htbp]
\centering
\caption{选择性终端层 Fresh-D2D 交叉实例验证（train D2D=0.02, seed=42）。}
\label{tab:work2_fresh_d2d}
\begin{tabular}{@{}lcccccc@{}}
\toprule
配置 & 评估 D2D & $n$ 种子 & 均值 PPL & 标准差 & min & max \\
\midrule
last1 [23] & 0.02 & 5 & 18.42 & 0.02 & 18.40 & 18.45 \\
last1 [23] & 0.04 & 5 & 18.55 & 0.03 & 18.52 & 18.58 \\
last1 [23] & 0.05 & 5 & 18.60 & 0.03 & 18.56 & 18.63 \\
\midrule
last2 [22,23] & 0.02 & 5 & 18.71 & 0.02 & 18.69 & 18.74 \\
last2 [22,23] & 0.04 & 5 & 19.07 & 0.04 & 19.03 & 19.13 \\
last2 [22,23] & 0.05 & 5 & 19.21 & 0.04 & 19.17 & 19.26 \\
\midrule
全层 & 0.02 & 5 & 26.05 & 0.53 & 25.17 & 26.50 \\
全层 & 0.04 & 5 & 44.34 & 2.65 & 40.08 & 46.74 \\
全层 & 0.05 & 5 & 66.84 & 5.91 & 57.65 & 73.06 \\
\bottomrule
\end{tabular}
\end{table}

这些结果阶段性支持假设~\ref{hyp:precision} 在选择性终端层范围内成立：当 $\sigma_{\text{d2d}}\le0.05$ 且量化为 8-bit 时，last1/last2 的困惑度退化处于可继续优化的范围，并且 fresh-D2D 方差很低。假设~\ref{hyp:retention} 和假设~\ref{hyp:endurance} 尚需时间动态和刷新实验确认，因此器件物理匹配目前应被表述为有先验支撑的候选优势，而不是最终部署结论。

\FloatBarrier

\section{贡献、边界与后续判据}
\label{sec:work2_contributions}
\label{sec:work2_boundaries}

Work~2 与 Work~1 的关系不是重复证明有机 CIM 可以运行神经网络，而是寻找更适合其物理特性的负载。Work~1 说明，动态激活路径、全局注意力和固定 D2D 实例会把有机阵列的非理想性放大为部署风险；Work~2 则转向静态 KV 缓存，利用其低频写入、高频读取和会话级生命周期。若选择性终端层路线在后续复现中继续成立，本文的知识增量在于说明同一种器件物理在不同负载下可从限制因素转化为候选优势。方法论增量在于把器件保持时间、耐久性、噪声和刷新机制映射到 LLM 困惑度、任务指标和能效，而不是只报告单器件指标或孤立的矩阵乘法误差。

当前边界见表~\ref{tab:boundaries}。这些排除项并不否定未来可扩展性，而是防止把行为级仿真和 1--3B 模型上的结论外推到尚未验证的系统规模。尤其是 $>$7B 模型、多片芯粒互连、128k 上下文和光电前端联合训练都可能成为后续工作，但它们需要新的容量模型、调度策略和硬件假设，不能由本章数据直接支撑。

\begin{table}[!htbp]
\centering
\caption{范围边界与排除理由。}
\label{tab:boundaries}
\small
\begin{tabularx}{\textwidth}{@{}l>{\raggedright\arraybackslash}X@{}}
\toprule
排除项 & 理由 \\
\midrule
$>$7B 参数模型 & 当前有机交叉阵列密度和仿真成本不足以支撑完整 KV 缓存评估；多片划分属于系统架构问题 \\
$>$32k token 上下文 & 超长上下文需要 3D 堆叠、混合存储层级或片外卸载，不再只考察有机 CIM 本征物理 \\
KV 缓存片上训练 & 训练需要高频原地更新 K/V 张量，超出 $10^3$--$10^4$ 次耐久假设 \\
生产部署声明 & 本章是行为级仿真和路线筛选，不给出商用部署承诺 \\
完整光电协同设计 & 光子前端与 CIM 权重联合训练属于更宽的前端方向，不作为本章独立贡献 \\
多片/芯粒扩展 & 芯粒互连和跨片调度属于工程外延，仅在架构探索中作为约束讨论 \\
\bottomrule
\end{tabularx}
\end{table}

后续判断标准也应保持可审计。当前选择性 KV 锁定包已完整归档 last1 至 all24 的逐行命令、检查点哈希、seed 语义和 JSON 文件，因此可用于支持层范围排序和全层负面对照结论；但时间动态实验仍需给出 $\tau$、刷新间隔、能耗和困惑度之间的曲线，才能继续支撑假设~\ref{hyp:retention} 和假设~\ref{hyp:endurance}。若这些实验不支持长期会话，本章结论应收缩为短会话或低层覆盖范围内的原型级可行性。

\FloatBarrier

\section{本章小结}
\label{sec:work2_summary}

本章把有机光电 CIM 从视觉权重映射扩展到 LLM KV 缓存场景。核心逻辑是：KV 缓存具有低频写入、高频读取和会话级生命周期，有机 OEC-RAM 的 $\sim10^3$~s 保持时间、有限耐久性和光学刷新机制与这一负载在数量级上更匹配。Remote~107 早期数据排除了全层模拟 KV 缓存作为主路线：8-bit 全层零噪声已达 PPL=17.48，相对旧数字基线 15.68 退化 11.5\%；6-bit 全层在 D2D/C2C 扫描中迅速失控。当前 selective-KV claim-lock packet 进一步锁定了相同排序：在 Pythia-410M、WikiText-2 raw-v1 test、context 512、non-overlapping stride 512、batch size 1 下，数字基线为 PPL=23.30；last1/last2 在 D2D=0.02--0.05 下保持在约 18.6--19.0 PPL，last4 升至约 20.4--21.5 PPL，而 all24 在 D2D=0.05 时恶化到 $59.055\pm1.610$ PPL。上述证据支持选择性终端层作为 Work~2 的主候选，并将全层模拟 KV 缓存定位为负面对照；第~\ref{chap:deployment-envelope}~章将回到视觉主线的部署包络，并把 Work~2 当前结论作为有机 CIM 适用性边界的一部分。
